\documentclass[10pt,letterpaper]{article}
\usepackage[utf8]{inputenc}
\usepackage[intlimits]{amsmath}
\usepackage{amsfonts}
\usepackage{amssymb}
\usepackage{ragged2e}
\usepackage[letterpaper, margin=0.5in]{geometry}
\usepackage{graphicx}
\usepackage{cancel}
\usepackage{mathtools}
\usepackage{tabularx}
\usepackage{arydshln}
\usepackage{tensor}
\usepackage{array}
\usepackage{xcolor}
\usepackage[boxed]{algorithm}
\usepackage[noend]{algpseudocode}
\usepackage{listings}
\usepackage{textcomp}
% \usepackage[pdf,tmpdir,singlefile]{graphviz}
\usepackage{mathrsfs}
\usepackage{bbm}
\usepackage{tikz}
\usepackage{tikz-cd}
\usepackage{enumitem}
\usepackage{arydshln}
\usepackage{relsize}
\usepackage{multicol}
\usepackage{scalerel}
\usepackage{upgreek}
% \usepackage{wasysym}
\usepackage{ifthen}

\usetikzlibrary{bayesnet}
\setlist{noitemsep}

%%%%%%%%%%%%%%%%%%%%%%%%%%%%%
% Formatting commands
%%%%%%%%%%%%%%%%%%%%%%%%%%%%%
\newcommand{\n}{\hfill\break}
\newcommand{\nn}{\vspace{0.5\baselineskip}\n}
\newcommand{\up}{\vspace{-\baselineskip}}
\newcommand{\hangblock}[2]{\par\noindent\settowidth{\hangindent}{\textbf{#1: }}\textbf{#1: }\nolinebreak#2}
\newcommand{\lemma}[1]{\hangblock{Lemma}{#1}}
\newcommand{\defn}[1]{\hangblock{Defn}{#1}}
\newcommand{\thm}[1]{\hangblock{Thm}{#1}}
\newcommand{\cor}[1]{\hangblock{Cor}{#1}}
\newcommand{\prop}[1]{\hangblock{Prop}{#1}}
\newcommand{\ex}[1]{\hangblock{Ex}{#1}}
\newcommand{\exer}[1]{\hangblock{Exer}{#1}}
\newcommand{\fact}[1]{\hangblock{Fact}{#1}}
\newcommand{\remark}[1]{\hangblock{Remark}{#1}}
\newcommand{\proven}{\;$\square$\n}
\newcommand{\problem}[1]{\par\noindent{\nolinebreak#1}\n}
\newcommand{\problempart}[2]{\par\noindent\indent{}\settowidth{\hangindent}{\textbf{(#1)} \indent{}}\textbf{(#1) }\nolinebreak#2\n}
\newcommand{\ptxt}[1]{\textrm{\textnormal{#1}}}
\newcommand{\inlineeq}[1]{\centerline{$\displaystyle #1$}}
\newcommand{\pageline}{\up\par\noindent\rule{\textwidth}{0.1pt}}

%%%%%%%%%%%%%%%%%%%%%%%%%%%%%
% Math commands
%%%%%%%%%%%%%%%%%%%%%%%%%%%%%
% Set Theory
\newcommand{\card}[1]{\left|#1\right|}
\newcommand{\set}[1]{\left\{#1\right\}}
\newcommand{\setmid}{\;\middle|\;}
\newcommand{\ps}[1]{\mathcal{P}\left(#1\right)}
\newcommand{\pfinite}[1]{\mathcal{P}^{\ptxt{finite}}\left(#1\right)}
\newcommand{\naturals}{\mathbb{N}}
\newcommand{\N}{\naturals}
\newcommand{\integers}{\mathbb{Z}}
\newcommand{\Z}{\integers}
\newcommand{\rationals}{\mathbb{Q}}
\newcommand{\Q}{\rationals}
\newcommand{\reals}{\mathbb{R}}
\newcommand{\R}{\reals}
\newcommand{\complex}{\mathbb{C}}
\newcommand{\C}{\complex}
\newcommand{\halfPlane}{\mathbb{H}}
\let\HSym\H
\let\H\relax
\newcommand{\H}{\halfPlane}
\newcommand{\comp}{^{\complement}}
\DeclareMathOperator{\Hom}{Hom}
\newcommand{\Ind}{\mathbbm{1}}
\newcommand{\cut}{\setminus}
\DeclareMathOperator{\elem}{elem}

% Differentiable Manifolds
\newcommand{\RP}{\mathbb{RP}}
\newcommand{\CP}{\mathbb{RP}}
\newcommand{\osubset}{\overset{\mathclap{\scalebox{0.5}{\ptxt{open}}}}{\subset}}
\newcommand{\osubseteq}{\overset{\mathclap{\scalebox{0.5}{\ptxt{open}}}}{\subseteq}}
\newcommand{\osupset}{\overset{\mathclap{\scalebox{0.5}{\ptxt{open}}}}{\supset}}
\newcommand{\osupseteq}{\overset{\mathclap{\scalebox{0.5}{\ptxt{open}}}}{\supseteq}}
\newcommand{\pdat}[3]{\left.\pd{#1}{#2}\right|_{#3}}
\DeclareMathOperator{\so}{so}
\DeclareMathOperator{\codim}{codim}
\DeclareMathOperator{\Diff}{Diff}
\let\dSym\d
\let\d\relax
\newcommand{\d}{\partial}
\DeclareMathOperator{\gl}{gl}

% Graph Theory
\let\deg\relax
\DeclareMathOperator{\deg}{deg}
\newcommand{\degp}{\ptxt{deg}^{+}}
\newcommand{\degn}{\ptxt{deg}^{-}}
\newcommand{\precdot}{\mathrel{\ooalign{$\prec$\cr\hidewidth\hbox{$\cdot\mkern0.5mu$}\cr}}}
\newcommand{\succdot}{\mathrel{\ooalign{$\cdot\mkern0.5mu$\cr\hidewidth\hbox{$\succ$}\cr\phantom{$\succ$}}}}
\DeclareMathOperator{\cl}{cl}
\DeclareMathOperator{\affdim}{affdim}

% Probability
\newcommand{\parSymbol}{\P}
\newcommand{\Prob}{\mathbb{P}}
\renewcommand{\P}{\Prob}
\newcommand{\Avg}{\mathbb{E}}
\newcommand{\E}{\Avg}
\DeclareMathOperator{\Var}{Var}
\DeclareMathOperator{\cov}{cov}
\DeclareMathOperator{\Unif}{Unif}
\DeclareMathOperator{\Binom}{Binom}
\newcommand{\CI}{\mathrel{\text{\scalebox{1.07}{$\perp\mkern-10mu\perp$}}}}

% Standard Math
\newcommand{\inv}{^{-1}}
\newcommand{\abs}[1]{\left|#1\right|}
\newcommand{\ceil}[1]{\left\lceil{}#1\right\rceil{}}
\newcommand{\floor}[1]{\left\lfloor{}#1\right\rfloor{}}
\newcommand{\conj}[1]{\overline{#1}}
\newcommand{\of}{\circ}
\newcommand{\tri}{\triangle}
\newcommand{\inj}{\hookrightarrow}
\newcommand{\surj}{\twoheadrightarrow}
\newcommand{\ndiv}{\nmid}
\renewcommand{\epsilon}{\varepsilon}
\newcommand{\divides}{\mid}
\newcommand{\ndivides}{\nmid}
\DeclareMathOperator{\lcm}{lcm}
\DeclareMathOperator{\sgn}{sgn}
\newcommand{\map}[4]{\!\!\!\begin{array}[t]{rcl}#1 & \!\!\!\!\to & \!\!\!\!#2\\ {}#3 & \!\!\!\!\mapsto & \!\!\!\!#4\end{array}}
\newcommand{\bigsum}[2]{\smashoperator[lr]{\sum_{\scalebox{#1}{$#2$}}}}
\DeclareMathOperator{\gcf}{gcf}
\newcommand{\restr}[2]{\left.#1\right|_{#2}}

% Linear Algebra
\newcommand{\Id}{\textrm{\textnormal{Id}}}
\newcommand{\im}{\textrm{\textnormal{im}}}
\newcommand{\norm}[1]{\abs{\abs{#1}}}
\newcommand{\tpose}{^{T}\!}
\newcommand{\iprod}[1]{\left<#1\right>}
\newcommand{\giprod}{\iprod{\;\,,\;}}
\DeclareMathOperator{\tr}{tr}
\DeclareMathOperator{\trace}{tr}
\newcommand{\chgBasMat}[3]{\!\!\tensor*[_{#1}]{\left[#2\right]}{_{#3}}}
\newcommand{\vecBas}[2]{\tensor*[]{\left[#1\right]}{_{#2}}}
\DeclareMathOperator{\GL}{GL}
\DeclareMathOperator{\Mat}{Mat}
\DeclareMathOperator{\vspan}{span}
\DeclareMathOperator{\rank}{rank}
\newcommand{\V}[1]{\vec{#1}}
\DeclareMathOperator{\proj}{proj}
\DeclareMathOperator{\compProj}{comp}
\DeclareMathOperator{\row}{row}
\newcommand{\smallPMatrix}[1]{\paren{\begin{smallmatrix}#1\end{smallmatrix}}}
\newcommand{\smallBMatrix}[1]{\brack{\begin{smallmatrix}#1\end{smallmatrix}}}
\newcommand{\pmat}[1]{\begin{pmatrix}#1\end{pmatrix}}
\newcommand{\bmat}[1]{\begin{bmatrix}#1\end{bmatrix}}
\newcommand{\dual}{^{*}}
\newcommand{\pinv}{^{\dagger}}
\newcommand{\horizontalMatrixLine}{\ptxt{\rotatebox[origin=c]{-90}{$|$}}}
\DeclareMathOperator{\range}{range}
\DeclareMathOperator{\Symm}{Symm}
\DeclareMathOperator{\SU}{SU}

% Multilinear Algebra
\let\Lsym\L
\let\L\relax
\DeclareMathOperator{\L}{\mathscr{L}}
\DeclareMathOperator{\A}{\mathcal{A}}
\DeclareMathOperator{\Alt}{Alt}
\DeclareMathOperator{\Sym}{Sym}
\newcommand{\ot}{\otimes}
\newcommand{\ox}{\otimes}
\DeclareMathOperator{\asc}{asc}
\DeclareMathOperator{\asSet}{set}
\DeclareMathOperator{\sort}{sort}
\DeclareMathOperator{\ringA}{\mathring{A}}
\DeclareMathOperator{\Sh}{Sh}

% Topology
\newcommand{\closure}[1]{\overline{#1}}
\newcommand{\uball}{\mathcal{U}}
\DeclareMathOperator{\Int}{Int}
\DeclareMathOperator{\Ext}{Ext}
\DeclareMathOperator{\Bd}{Bd}
\DeclareMathOperator{\rInt}{rInt}
\DeclareMathOperator{\ch}{ch}
\DeclareMathOperator{\ah}{ah}
\newcommand{\LargerTau}{\mathlarger{\mathlarger{\mathlarger{\mathlarger{\tau}}}}}
\newcommand{\Tau}{\mathcal{T}}

% Analysis
\DeclareMathOperator{\Graph}{Graph}
\DeclareMathOperator{\epi}{epi}
\DeclareMathOperator{\hypo}{hypo}
\DeclareMathOperator{\supp}{supp}
\newcommand{\lint}[2]{\underset{#1}{\overset{#2}{{\color{black}\underline{{\color{white}\overline{{\color{black}\int}}\color{black}}}}}}}
\newcommand{\uint}[2]{\underset{#1}{\overset{#2}{{\color{white}\underline{{\color{black}\overline{{\color{black}\int}}\color{black}}}}}}}
\newcommand{\alignint}[2]{\underset{#1}{\overset{#2}{{\color{white}\underline{{\color{white}\overline{{\color{black}\int}}\color{black}}}}}}}
\newcommand{\extint}{\ptxt{ext}\int}
\newcommand{\extalignint}[2]{\ptxt{ext}\alignint{#1}{#2}}
\newcommand{\conv}{\ast}
\newcommand{\pd}[2]{\frac{\partial{}#1}{\partial{}#2}}
\newcommand{\del}{\nabla}
\DeclareMathOperator{\grad}{grad}
\DeclareMathOperator{\curl}{curl}
\let\div\relax
\DeclareMathOperator{\div}{div}
\DeclareMathOperator{\vol}{vol}

% Complex Analysis
\let\Re\relax
\DeclareMathOperator{\Re}{Re}
\let\Im\relax
\DeclareMathOperator{\Im}{Im}
\DeclareMathOperator{\Res}{Res}

% Abstract Algebra
\DeclareMathOperator{\ord}{ord}
\newcommand{\generated}[1]{\left<#1\right>}
\newcommand{\cycle}[1]{\smallPMatrix{#1}}
\newcommand{\id}{\ptxt{id}}
\newcommand{\iso}{\cong}
\DeclareMathOperator{\Aut}{Aut}
\DeclareMathOperator{\SL}{SL}
\DeclareMathOperator{\op}{op}
\newcommand{\isom}[4]{\!\!\!\begin{array}[t]{rcl}#1 & \!\!\!\!\overset{\sim}{\to} & \!\!\!\!#2\\ #3 & \!\!\!\!\mapsto & \!\!\!\!#4\end{array}}
\newcommand{\F}{\mathbb{F}}
\newcommand{\acts}{\;\reflectbox{\rotatebox[origin=c]{-90}{$\circlearrowright$}}\;}
\newcommand{\disjunion}{\mathrel{\text{\scalebox{1.07}{$\perp\mkern-10mu\perp$}}}}
\DeclareMathOperator{\SO}{SO}
\DeclareMathOperator{\stab}{stab}
\DeclareMathOperator{\nullity}{nullity}
\DeclareMathOperator{\Perm}{Perm}
\DeclareMathOperator{\nsubgp}{\vartriangleleft}
\DeclareMathOperator{\notnsubgp}{\ntriangleleft}
\newcommand{\presentation}[2]{\left<#1\;\middle|\;#2\right>}
\DeclareMathOperator{\Char}{char}
\DeclareMathOperator{\fchar}{char}
\DeclareMathOperator{\triv}{triv}
\DeclareMathOperator{\reg}{reg}
\DeclareMathOperator{\std}{std}
\DeclareMathOperator{\Func}{Func}
\DeclareMathOperator{\End}{End}

% Convex Optimization
\let\sectionSymbol\S
\let\S\relax
\newcommand{\S}{\mathbb{S}}
\DeclareMathOperator{\dist}{dist}
\DeclareMathOperator{\dom}{dom}
\DeclareMathOperator{\diag}{diag}
\DeclareMathOperator{\ones}{\mathbbm{1}}
\newcommand{\minimizeOver}[3]{\begin{array}{rl}\underset{#1}{\ptxt{minimize}} & #2\\ \ptxt{subject to} & #3\end{array}}
\newcommand{\maximizeOver}[3]{\begin{array}{rl}\underset{#1}{\ptxt{maximize}} & #2\\ \ptxt{subject to} & #3\end{array}}
\newcommand{\minimizationProblem}[2]{\minimizeOver{}{#1}{#2}}
\newcommand{\maximizationProblem}[2]{\maximizeOver{}{#1}{#2}}
\newcommand{\minimizeOverUnconstrained}[2]{\begin{array}{rl}\underset{#1}{\ptxt{minimize}} & #2\end{array}}
\newcommand{\maximizeOverUnconstrained}[2]{\begin{array}{rl}\underset{#1}{\ptxt{maximize}} & #2\end{array}}
\newcommand{\minimizationUnconstrained}[1]{\minimizeOverUnconstrained{}{#1}}
\newcommand{\maximizationUnconstrained}[1]{\maximizeOverUnconstrained{}{#1}}
\DeclareMathOperator{\argmin}{argmin}
\DeclareMathOperator{\argmax}{argmax}

% Proofs
\newcommand{\st}{s.t.}
\newcommand{\unique}{!}
\newcommand{\iffdef}{\overset{\ptxt{def}}{\Leftrightarrow}}
\newcommand{\eqVertical}{\rotatebox[origin=c]{90}{=}}
\newcommand{\mapsfrom}{\mathrel{\reflectbox{\ensuremath{\mapsto}}}}
\newcommand{\mapsdown}{\rotatebox[origin=c]{-90}{$\mapsto$}\mkern2mu}
\newcommand{\mapsup}{\rotatebox[origin=c]{90}{$\mapsto$}\mkern2mu}
\newcommand{\from}{\!\mathrel{\reflectbox{\ensuremath{\to}}}}
\newcommand{\labeledeq}[1]{\overset{\mathclap{\ptxt{#1}}}{=}}
\newcommand{\eqdef}{\labeledeq{def}}

% Brackets
\newcommand{\paren}[1]{\left(#1\right)}
\renewcommand{\brack}[1]{\left[#1\right]}
\renewcommand{\brace}[1]{\left\{#1\right\}}
\newcommand{\ang}[1]{\left<#1\right>}

% Algorithms
\algrenewcommand{\algorithmiccomment}[1]{\hskip 1em \texttt{// #1}}
\algrenewcommand\algorithmicrequire{\textbf{Input:}}
\algrenewcommand\algorithmicensure{\textbf{Output:}}
\newcommand{\algP}{\ptxt{\textbf{P}}}
\newcommand{\algNP}{\ptxt{\textbf{NP}}}
\newcommand{\algNPC}{\ptxt{\textbf{NP-Complete}}}
\newcommand{\algNPH}{\ptxt{\textbf{NP-Hard}}}
\newcommand{\algEXP}{\ptxt{\textbf{EXP}}}
\DeclareMathOperator{\fl}{fl}

%%%%%%%%%%%%%%%%%%%%%%%%%%%%%
% Other commands
%%%%%%%%%%%%%%%%%%%%%%%%%%%%%
\newcommand{\flag}[1]{\textbf{\textcolor{red}{#1}}}
\let\uSym\u
\let\u\relax
\newcommand{\u}[1]{\underline{#1}}
\let\bSym\b
\let\b\relax
\newcommand{\b}[1]{\textbf{#1}}
\let\iSym\i
\let\i\relax
\newcommand{\i}[1]{\textit{#1}}
\let\scSym\sc
\let\sc\relax
\newcommand{\sc}[1]{\textsc{#1}}

%%%%%%%%%%%%%%%%%%%%%%%%%%%%%%%%%%%%%%%
% Make l's curvy in math environments %
%%%%%%%%%%%%%%%%%%%%%%%%%%%%%%%%%%%%%%%
\mathcode`l="8000
\begingroup
\makeatletter
\lccode`\~=`\l
\DeclareMathSymbol{\lsb@l}{\mathalpha}{letters}{`l}
\lowercase{\gdef~{\ifnum\the\mathgroup=\m@ne \ell \else \lsb@l \fi}}%
\endgroup

%%%%%%%%%%%%%%%%%%%%%%%%%
% Fix \vdots and \ddots %
%%%%%%%%%%%%%%%%%%%%%%%%%
\usepackage{letltxmacro}
\LetLtxMacro\orgvdots\vdots
\LetLtxMacro\orgddots\ddots

\makeatletter
\DeclareRobustCommand\vdots{%
	\mathpalette\@vdots{}%
}
\newcommand*{\@vdots}[2]{%
	% #1: math style
	% #2: unused
	\sbox0{$#1\cdotp\cdotp\cdotp\m@th$}%
	\sbox2{$#1.\m@th$}%
	\vbox{%
		\dimen@=\wd0 %
		\advance\dimen@ -3\ht2 %
		\kern.5\dimen@
		% remove side bearings
		\dimen@=\wd2 %
		\advance\dimen@ -\ht2 %
		\dimen2=\wd0 %
		\advance\dimen2 -\dimen@
		\vbox to \dimen2{%
			\offinterlineskip
			\copy2 \vfill\copy2 \vfill\copy2 %
		}%
	}%
}
\DeclareRobustCommand\ddots{%
	\mathinner{%
		\mathpalette\@ddots{}%
		\mkern\thinmuskip
	}%
}
\newcommand*{\@ddots}[2]{%
	% #1: math style
	% #2: unused
	\sbox0{$#1\cdotp\cdotp\cdotp\m@th$}%
	\sbox2{$#1.\m@th$}%
	\vbox{%
		\dimen@=\wd0 %
		\advance\dimen@ -3\ht2 %
		\kern.5\dimen@
		% remove side bearings
		\dimen@=\wd2 %
		\advance\dimen@ -\ht2 %
		\dimen2=\wd0 %
		\advance\dimen2 -\dimen@
		\vbox to \dimen2{%
			\offinterlineskip
			\hbox{$#1\mathpunct{.}\m@th$}%
			\vfill
			\hbox{$#1\mathpunct{\kern\wd2}\mathpunct{.}\m@th$}%
			\vfill
			\hbox{$#1\mathpunct{\kern\wd2}\mathpunct{\kern\wd2}\mathpunct{.}\m@th$}%
		}%
	}%
}
\makeatother

\newcommand{\B}{
	\begin{tikzpicture}
	\filldraw [fill=red, draw=black] (0, 0) rectangle (0.37, 0.45);
	\draw [line width=0.5mm, white ] (0.1,0.08) -- (0.1,0.38);
	\draw[line width=0.5mm, white ] (0.1, 0.35) .. controls (0.2, 0.35) and (0.4, 0.2625) .. (0.1, 0.225);
	\draw[line width=0.5mm, white ] (0.1, 0.225) .. controls (0.2, 0.225) and (0.4, 0.1625) .. (0.1, 0.1);
	\end{tikzpicture}
}

% Allow custom symbols as arrows in tikz-cd
\tikzset{
  symbol/.style={
    draw=none,
    every to/.append style={
      edge node={node [sloped, allow upside down, auto=false]{$#1$}}}
  }
}
% Example: \arrow[r,symbol=\cong]
% https://tex.stackexchange.com/questions/394154/how-to-include-inclusion-subgroup-relationship-in-tikz-cd-diagram

\author{Thomas Cohn}
\title{Math 591 Lecture 37}
\date{12/4/20} % Can also use \today

\begin{document}
\maketitle
\setlength\RaggedRightParindent{\parindent}
\RaggedRight

\par\noindent
Briefly, think back to integration of forms. Observe:
\[
	\int_{M}\mu=\sum_{\alpha}\int(\phi_{\alpha}\inv)^{*}(\chi_{\alpha}\mu)
\]
We claimed (without proof) that this sum has finitely-many nonzero summands.\nn
Proof: Let $K=\supp\mu$ (note that $K$ must be compact), and $U_{\alpha}=\supp\chi_{\alpha}$. Note that $\set{U_{\alpha}}$ is locally finite: $\forall{}p\in{}K$, $\exists{}V_{p}$, a neighborhood of $p$, such that $\set{\alpha\mid{}V_{p}\cap{}U_{\alpha}\ne\emptyset}$ is finite. Also, $\set{V_{p}}_{p\in{}K}$ is naturally a cover of $K$, so by compactnes, $\exists{}p_{1},\ldots,p_{n}\in{}K$ \st{} $\set{V_{p_{1}},\ldots,V_{p_{n}}}$ is a cover of $K$. Finally,
\[
	\set{\alpha\mid{}K\cap{}U_{\alpha}\ne\emptyset}\subset\bigcup_{j=1}^{n}\underbrace{\set{\alpha\mid{}V_{p_{j}}\cap{}U_{\alpha}\ne\emptyset}}_{\ptxt{finite}}
\]
\proven

\par\noindent
Now, we return to working with manifolds with boundary...\n

\par\noindent
The principle is that all definitions on manifold with boundary are exact analogues to the case where $\d{}M=\emptyset$. The key difference is the model spaces are open subsets of $\H^{n}=\set{x^{n}\ge{}0}\subseteq\R^{n}$, and transition maps are diffeomorphisms between open subsets of $\H^{n}$ (which carry boundary points to boundary points, and restrict to diffeomorphisms between open sets of $\d\H^{n}\cong\R^{n-1}$).\n

\ex{
	Notions of smoothness. $f:M\to\R$ is smooth iff for any chart $\phi_{\alpha}$, $f_{\alpha}=f\of\phi_{\alpha}\inv:U\to\R$ is smooth (with $U\osubseteq\H^{n}$). In case $U\cap\d\H^{n}\ne\emptyset$, smoothness of $f_{\alpha}:U\to\R$ means $\exists\tilde{f}_{\alpha}:\tilde{U}_{\alpha}\to\R$, a smooth extension of $f_{\alpha}$ to $\tilde{U}_{\alpha}\osubseteq\R^{n}$. Then $U=\tilde{U}_{\alpha}\cap\H^{n}$.\n
}

\ex{
	Existence of partitions of unity, just as before.\n
}

\section*{Tangent Spaces}

\par\noindent
$\forall{}p\in\d{}M$, $T_{p}M$ is still $n$-dimensional, and is still spanned by $\set{\pd{}{x^{1}},\ldots,\pd{}{x^{n}}}$. $\pd{}{x^{n}}$ is well-defined on smooth functions, because such functions extend across the boundary, and last time, we showed the choice of extension doesn't matter.\n

\par\noindent
Note: $\d{}M$ inherits a $C^{\infty}$ manifold structure, by restricting charts of $M$.\n

\par\noindent
$\forall{}p\in\d{}M$, $T_{p}(\d{}M)\overset{\iota_{*,p}}{\hookrightarrow}T_{p}M$, where $\iota:\d{}M\to{}M$ is the inclusion. Identify $T_{p}(\d{}M)\subset{}T_{p}M$ as a subspace -- the image of $\iota_{*,p}$. It's a hyperplane, i.e., $T_{p}(\d{}M)$ has codimension $1$.\n

\par\noindent
Note: $(T_{p}M)\cut(T_{p}\d{}M)$ has two components: ``inward-pointing'' vectors and ``outward-pointing'' vectors.\n

\par\noindent
How do we characterize outward-pointing vectors? Well, say $p\in\d{}M$. Let $\gamma:(-\varepsilon,0]\to{}M$ smooth with $\gamma(0)=p$. Then define $\dot\gamma(0)$ to be outward pointing.\n

\subsection*{Orientation of Manifolds with Boundary}

\par\noindent
Once again, this is an exact analogue of manifolds without boundary. $\forall{}p\in{}M$, $T_{p}M$ has dimension $n$, two orientations, etc.\n

\lemma{
	If $M$ is orientable, then $\d{}M$ is orientable.\n
}

\par\noindent
By convention, $\forall{}p\in\d{}M$, a basis $\paren{b_{1},\ldots,b_{n-1}}$ of $T_{p}(\d{}M)$ is positive iff $\paren{\nu,b_{1},\ldots,b_{n-1}}$ is a positive basis of $T_{p}M$, \i{for any outward-pointing $\nu\in{}T_{p}M$}. This defines the boundary's orientation.\n

\ex{
	Say $M=\H^{n}$, oriented such that $\set{\pd{}{x^{1}},\ldots,\pd{}{x^{n}}}$ is positive. What is the boundary orientation?\nn
	For $n=2$, $\set{\pd{}{x^{1}},\pd{}{x^{2}}}$ is positive. So if $\nu$ is outward-pointing, $\set{\nu,\pd{}{x^{2}}}$ is positive, so $\set{\pd{}{x^{1}}}$ is positive on $\d\H^{2}\cong\R$.\n
	For $n=3$, is $(\vec{i},\vec{j})$ positive? Look at $v=-\vec{k}$ -- $(-\vec{k},\vec{i},\vec{j})$ is negative. So we claim that $\d\H^{n}$ is $(-1)^{n}$ times the standard orientation of $\R^{n}$.\n
}

\section*{Stokes' Theorem}

\thm{
	(Stokes' Theorem) Let $M$ be a manifold (possibly with boundary). Let $\mu\in\Omega_{0}^{n-1}(M)$. Assume $M$ is oriented; give $\d{}M$ the boundary orientation. Then
	\[
		\int_{\d{}M}\iota^{*}\mu=\int_{M}d\mu
	\]
	(where $\iota:\d{}M\hookrightarrow{}M$ is the inclusion map). One often omits the $\iota^{*}$, so we say
	\[
		\int_{\d{}M}\mu=\int_{M}d\mu
	\]
	(If $\d{}M=\emptyset$, then $\int_{M}d\mu=0$.)\n
}

\par\noindent
Observe: $\iota^{*}\mu$ is a top degree form on the boundary.\n

\ex{
	$M=[a,b]$, $n=1$, $\mu\in\Omega_{0}^{0}(M)\cong{}C^{\infty}(M)$. Then $\int_{\d{}M}\mu=\mu(b)-\mu(a)=\int_{M}df$. The ``$-$'' sign comes because of the orientation of $\d{}M$ is outward-pointing.\n
}

\ex{
	$n=2$. Then we get Green's Theorem:\n
	$\mu=P\,dx+Q\,dy$. $d\mu=(Q_{x}-P_{y})\,dx\wedge{}dy$.\n
	The orientation on $\d{}M$ comes from the right-hand rule.
	\[
		\int_{\d{}M}P\,dx+Q\,dy=\iint_{M}(Q_{x}-P_{y})\,dx\wedge{}dy
	\]
	\n
}

\exer{
	For $n=3$, we get the usual Stokes' theorem.\n
}

\subsection*{Proof of Stokes' Theorem}

\par\noindent
First, assume it holds for $\H^{n}$. Then it follows for $\mu\in\Omega_{0}^{n-1}(M)$ if $\mu$ is supported in the domain of a chart $\phi$:
\[
	\int_{\d{}M}\mu=\int_{\d\H^{n}}(\phi\inv)^{*}\iota^{*}\mu\labeledeq{(1)}\int_{\H^{n}}d(\phi\inv)^{*}\mu=\int_{\H^{n}}(\phi\inv)^{*}d\mu=\int_{M}d\mu
\]
with (1) by the definition of integrals, plus our assumption about $\mu$. In general (still under the assumption that Stokes' theorem holds for $\H^{n}$), we use a partition of unity $\set{\chi_{j}}$ subordinate to an atlas.
\[
	\int_{\d{}M}\iota^{*}\mu=\sum_{j}\int_{\d{}M}\iota^{*}(\chi_{j}\mu)=\sum_{j}\int_{M}d(\chi_{j}\mu)=\sum_{j}\int_{M}d(\chi_{j})\wedge\omega+\int_{M}\chi_{j}d\mu=\int_{M}d\mu+\int_{m}\left(\vphantom{\sum_{j}}\right.\underbrace{\sum_{j}d\chi_{j}}_{\substack{=d\sum_{j}\chi_{j}=0\\ \ptxt{b/c }\sum_{j}\chi_{j}\equiv{}1}}\left.\vphantom{\sum_{j}}\right)\wedge\mu=\int_{M}d\mu
\]
So finally, it's enough to prove Stokes' Theorem for $M=\H^{n}$. Write
\[
	\mu=\sum_{i=1}^{n}a_{i}dx^{1}\wedge\cdots\wedge\widehat{dx^{i}}\wedge\cdots\wedge{}dx^{n},\quad{}a_{i}\in{}C^{\infty}(\H^{n})
\]
Then
\[
	d\mu=\paren{\sum_{i=1}^{n}\pd{a_{i}}{x^{j}}(-1)^{j-1}}\underbrace{dx^{1}\wedge\cdots\wedge{}dx^{n}}_{\ptxt{standard volume form}}
\]
Because we assume $\mu$ has compact support, $\supp\mu$ is bounded, so $\exists{}R>0$ \st{}
\[
	\supp\mu\subseteq\underbrace{[-R,R]\times\cdots\times[-R,R]}_{n-1\ptxt{ times}}\times[0,R]
\]
Thus,
\[
	\int_{\H^{n}}d\mu=\sum_{i=1}^{n}(-1)^{i-1}\int_{0}^{R}\int_{-R}^{R}\cdots\int_{-R}^{R}\pd{a_{i}}{x^{j}}\,dx^{1}\cdots{}dx^{n}
\]
So $\forall{}i$, first do $\int_{-R}^{R}\pd{a_{i}}{x^{i}}\,dx^{i}$. For $i=1,\ldots,n-1$,
\[
	\int_{-R}^{R}\pd{a_{i}}{x^{i}}\,dx^{i}=a_{}(x^{1},\ldots,R,\ldots,x^{n})-a_{i}(x^{1},\ldots,-R,\ldots,x^{n})=0-0=0
	\qquad
	\int_{0}^{R}\pd{a_{n}}{x^{n}}\,dx^{n}=-a_{n}(x^{1},\ldots,x^{n-1},R)
\]
Thus,
\[
	\int_{\H^{n}}d\mu=\underbrace{-(-1)^{n-1}}_{=(-1)^{n}}\int_{\R^{n-1}}a_{n}(x^{1},\ldots,x^{n-1},0)\,dx^{1}\cdots{}dx^{n-1}
\]
And on the other hand,
\[
	\int_{\d\H^{n}}\mu=\underbrace{(-1)^{n}}_{\mathclap{\substack{\uparrow\\ \ptxt{boundary orientation}}}}\int_{\R^{n-1}}a_{n}(x^{1},\ldots,x^{n-1},0)\,dx^{1}\cdots{}dx^{n-1}
\]
(with the $0$ appearing in the $a_{n}$ term because $\iota^{*}(dx^{n})=0$.)\n
\proven

\end{document}